% Fonts
\usepackage[utf8]{inputenc}
\usepackage[spanish,es-tabla,es-nodecimaldot]{babel}
\usepackage{mathpazo}% Palatino font
\usepackage[T1]{fontenc}
\usepackage[autostyle=true]{csquotes}

% Desing and tables
\usepackage{ragged2e}
\usepackage{float}
\usepackage{lastpage}
\usepackage{caption}
\usepackage{titlesec}
\usepackage{listings}
\usepackage{tabularx}
\usepackage[table,xcdraw]{xcolor}
\usepackage{pdflscape}
\usepackage{pdfpages}
\usepackage[left=2cm,top=2.5cm,right=2cm,bottom=2.5cm]{geometry}

% Maths
\usepackage{amsmath,amsthm,amsfonts,amssymb}
\usepackage{bm}



% Graphics and gantt
\usepackage{enumerate, graphicx,subcaption}
\usepackage{pgfgantt}

% Bibliography
\usepackage[
    backend=biber,      % Motor recomendado
    style=apa,          % Estilo APA 7ma edición
    language=spanish    % Para que ponga "&" y "Recuperado de" correctamente
]{biblatex}
\addbibresource{Referencias/referencias.bib}


% Hyperref
\usepackage{hyperref}

\hypersetup{
  colorlinks=true,
  linkcolor=blue,
  filecolor=magenta,
  urlcolor=blue,
  citecolor=blue,
  pdfauthor={\estudiante},
  pdftitle={\matricula\ - \estudiante - \curso},
  pdfsubject={\curso},
  pdfkeywords={\pclaves}
}

% Este código envuelve el comando de cita de APA en un hipervínculo completo
\makeatletter
\DeclareCiteCommand{\parencite}[\mkbibparens]
  {\usebibmacro{cite:init}%
   \usebibmacro{prenote}}
  {\usebibmacro{citeindex}%
   \printtext[bibhyperref]{\usebibmacro{cite}}}
  {\multicitedelim}
  {\usebibmacro{postnote}}

\DeclareCiteCommand{\cbx@textcite}
  {\usebibmacro{cite:init}}
  {\usebibmacro{citeindex}%
   \printtext[bibhyperref]{\usebibmacro{textcite}}}
  {\multicitedelim}
  {\usebibmacro{postnote}}
\makeatother


% Numerar con sección
\numberwithin{figure}{section}
\numberwithin{table}{section}
\numberwithin{equation}{section}
\setlength{\parskip}{0.3cm}
\setcounter{tocdepth}{4}
\setcounter{secnumdepth}{4}


% % Custom math commands
\newtheorem{definition}{Definición}
\newtheorem{theorem}{Teorema}
\newtheorem{lemma}{Lema}
\newtheorem{proposition}{Proposición}

% Índice de figuras/tablas como sección en vez de capítulo
\usepackage{tocbasic}
\addtotoclist[float]{lof}
\renewcommand*\listoffigures{\listoftoc[{\listfigurename}]{lof}}
\addtotoclist[float]{lot}
\renewcommand*\listoftables{\listoftoc[{\listtablename}]{lot}}
\setuptoc{lof}{leveldown}
\setuptoc{lot}{leveldown}

\usepackage{fancyhdr}
\pagestyle{fancy}
\fancyhf{}


\setlength{\headheight}{27.72116pt}
\addtolength{\topmargin}{-2.12115pt}


% pie de pagina corregido
\renewcommand{\footrulewidth}{0.5pt} % Línea divisoria en el pie
\fancyhf{}

\usepackage[font={footnotesize}]{caption}

% Custom fixes

\urlstyle{same}
\title{\titulo}
\author{\estudiante}
\date{\fecha}


\fancyfoot[L]{P\'agina \thepage \hspace{1pt} de \pageref{LastPage}}
\fancyfoot[C]{}
\fancyfoot[R]{}

% Si quieres que en las páginas pares/impares sea igual:
\lfoot{P\'agina \thepage \hspace{1pt} de \pageref{LastPage}}
\cfoot{}
\rfoot{}


\usepackage{xcolor}

% Definimos que el formato de los hipervínculos de cita abarque todo el bloque (#1)
\DeclareFieldFormat{citehyperref}{\bibhyperref{#1}}

% Creamos un alias para que biblatex use este formato en las citas
\DeclareFieldAlias{bibhyperref}{citehyperref}
